\documentclass[UTF8,10pt]{ctexart}
\usepackage{raremetals}

% 本文件是从期刊 Word 模板提炼的非官方 LaTeX 迁移入口。
% 投稿前请核对当期官方模板；编辑部字段不要由作者虚构。
\RMTitleCN{【中文题目：官网要求不超过20字】}
\RMAuthorCN{【作者姓名】}
\RMAffiliationCN{（【单位全称，城市 邮政编码】）}
\RMAbstractCN{【官网要求中文摘要超过350字，应说明研究问题、过程和方法、主要结果与结论，并填写真实定量结果。】}
\RMKeywordsCN{【关键词1；关键词2；关键词3】}% 官网要求3--6个
\RMClassification{【中图分类号】}

\RMTitleEN{English Title within 200 characters}
\RMAuthorEN{Author One, Author Two}
\RMAffiliationEN{(Affiliation, City Postal Code, China)}
\RMAbstractEN{Insert the 800--1000-word English long abstract required by the journal. Do not copy the Chinese abstract mechanically without checking the current instructions.}
\RMKeywordsEN{keyword one; keyword two; keyword three}

\RMFunding{【基金名称及编号；无则按期刊要求处理】}
\RMAuthorBio{【第一作者姓名、出生年、性别、生源地、学历或职称、研究方向和电子邮箱】}
\RMCorrespondingAuthor{【姓名、职称、电话和电子邮箱】}

\begin{document}
\makermfrontmatter

% 官网要求引言200--400字，不编号、不显示“引言”二字，也不叙述结果或结论。
【在这里写引言正文，并使用真实引用\cite{example2026}。】

\section{实验}
\subsection{实验材料}
【填写实验材料与方法。】

\begin{table}[htbp]
  \centering
  \bilingualcaption{中文表题}{English table title}
  \begin{tabular}{ccc}
    \toprule
    项目 & 数值 & 单位 \\
    \midrule
    【填写】 & 【填写】 & 【填写】 \\
    \bottomrule
  \end{tabular}
  \label{tab:example}
\end{table}

\section{结果与讨论}
\subsection{结果}
【填写结果与讨论；图表必须在正文中引用。】

\begin{figure}[htbp]
  \centering
  \fbox{\rule{0pt}{35mm}\rule{0.72\linewidth}{0pt}}
  \bilingualcaption{中文图题}{English figure title}
  \label{fig:example}
\end{figure}

\section{结论}
【填写由数据直接支持的结论。】

\begin{rmreferences}
\bibliographystyle{unsrt}
\bibliography{bib/references}
\end{rmreferences}

\makermenglish
\makermmetadata
\end{document}
